\chapter{Cybersécurité}

La cybersécurité est un axe central du projet, tant dans l'outil SRS lui-même que dans
les exigences qu'il génère pour la future plateforme.

\section{Modèle de menace}
Le moteur de sécurité de SRS applique une analyse de type \textbf{STRIDE} (Spoofing,
Tampering, Repudiation, Information disclosure, Denial of service, Elevation of
privilege), pondérée selon une logique \textbf{DREAD}, sur un catalogue de concepts
prioritaires du Knowledge Graph (authentification, autorisation, protection des
données, disponibilité, etc.).

\section{Analyse des risques}
Pour chaque risque identifié, un registre structuré est généré, incluant~: identifiant,
actif concerné, menace, vulnérabilité, impact, probabilité, niveau de risque,
mitigation proposée et risque résiduel. Seuls les risques dont le niveau est
\emph{élevé} ou \emph{critique} apparaissent dans le registre final, afin de ne pas
diluer l'attention sur des risques mineurs.

\section{Contrôle d'accès}
Une matrice RBAC (Role-Based Access Control) est générée à partir des acteurs
effectivement identifiés dans les exigences fonctionnelles — elle n'est donc pas fixe,
mais reflète les rôles confirmés durant l'entretien (par exemple~: administrateur,
coordinateur régional, représentant de coopérative).

\section{Protection des données}
Une analyse d'impact relative à la protection des données (PIA) est générée, avec une
règle stricte~: les exigences légales telles que la Loi~09-08 ou le RGPD ne sont
\textbf{jamais} présentées comme confirmées par défaut. Elles sont systématiquement
étiquetées comme « à confirmer avec le responsable juridique / DPO », sauf confirmation
explicite d'une partie prenante durant l'entretien.

\section{Sécurité des fichiers}
\TODO{Détailler les contrôles de sécurité relatifs au téléversement de fichiers
(validation de type, taille, scan antivirus) une fois ces exigences confirmées pour un
projet concret — ne pas décrire de contrôles qui ne sont pas encore documentés.}

\section{Journalisation et traçabilité}
Chaque exigence générée référence la réponse d'entretien dont elle est issue
(\texttt{source\_concept\_key}), fournissant une traçabilité de bout en bout entre une
exigence du Cahier des Charges et la réponse de la partie prenante qui l'a motivée.

\section{Sauvegarde et continuité}
\TODO{Documenter la stratégie de sauvegarde et de continuité effectivement retenue pour
l'environnement d'hébergement de SRS, une fois le déploiement finalisé.}

\section{Secure SDLC}
Le développement de SRS applique des pratiques de cycle de développement sécurisé de
base~: séparation stricte de la configuration sensible (clés API) via variables
d'environnement, absence de journalisation des secrets, et couverture par tests
automatisés des principaux flux applicatifs (voir chapitre~8).
